% 嵌套分组 SDUSetup 回归测试
% 验证 \SDUSetup{ info={...}, option={...} } 的嵌套写法能正确设置各变量。
% 同时验证平铺写法与嵌套写法等价（同一组内部变量）。
\documentclass{sduthesis}
\input{regression-test}

\SDUSetup{
  module = {undergraduate},
  info   = {
    title      = {嵌套分组论文},
    author     = {嵌套作者},
    studentId  = {202200000001},
    school     = {信息学院},
    major      = {计算机},
    supervisor = {王教授},
    year       = {2025},
    month      = {6},
  },
  option = {
    lineSpread = 1.5,
    pageLeft   = 3cm,
  },
}

\begin{document}
\START
% 嵌套 info 分组应正确写入各 Getter 变量
\ASSERT{\GetTitle}{嵌套分组论文}
\ASSERT{\GetAuthor}{嵌套作者}
\ASSERT{\GetStudentId}{202200000001}
\ASSERT{\GetSchool}{信息学院}
\ASSERT{\GetMajor}{计算机}
\ASSERT{\GetSupervisor}{王教授}
\ASSERT{\GetYear}{2025}
\ASSERT{\GetMonth}{6}

% 嵌套 option 分组应覆盖样式变量（行为验证：不报错即通过）
\TYPE{nested-option-ok=YES}

\makecoverpage
\END
\end{document}
